\section{Dynamic models}
\label{s:models-dynamic}

The previous section's static models requires a strict factor exogeneity assumption
outlined in equation \eqref{eq:static-identification-factor-exogeneity}
for all firms, including the own-firm effect of the production function.
In this section, I use a dynamic model to relax this assumption.
I also introduce a vector of temporal regressors:
an autoregressor of own-firm GVA as a component of its TFP,
and allowing for the past peer TFP components I analysed
in the static model to influence current-period TFP
in a lagged fashion.
The structure of this specification now resembles a vector panel,
which I will denote with VAR(p,q)\footnote{$q$-lags refers
to those on the peer spillover effect, rather than those of a moving average component
as is traditional of a VARMA(p,q) specification}.
\begin{align}
    \label{eq:dyn-main-reduced}
    z_{i,t} &= \alpha_i + \gamma_t + \beta_0 + \sum_{k=1}^p\phi_k y_{i,t-k}
        + \sum_{k=0}^q\mathbf z_{t-k}^{spill}\boldsymbol\theta_k + \varepsilon_{i,t}
    \\
    \label{eq:dyn-main-peer-effects}
    y_{i,t} &= \alpha_i + \gamma_t + \beta_0 + \beta_1 k_{i,t} + \beta_2 l_{i,t} 
                + \sum_{k=1}^p\phi_k y_{i,t-k} + \varepsilon_{i,t}
                \\ \notag &\qquad 
                + \sum_{k=0}^q \left[
                    \delta_k \sum_{j\neq i} w_{ij}y_{j,t}
                    + (\xi_k-\delta_k\beta_1) \cdot \sum_{j\neq i} w_{ij}k_{j,t}
                    + (\zeta_k-\delta_k\beta_2) \cdot \sum_{j\neq i} w_{ij}l_{j,t}
                \right]
\end{align}

Adding the autoregressor $y_{i,t}$ is also highly desirable.
The panel setup with firm fixed effects already accounts for own-firm factors,
and the presentation of $R^2$ including and excluding variation accounted for by the fixed effects
throughout this paper demonstrates this.
In the peer effects setting, autoregressive lags may also help control for reverse causality.
If $y_{i,t-1}$ is causal for $z_{jt}$,
or feasibly $z_{j,t-1}$ and $\mathbb E[z_{j,t-1}, z_{j,t}]\neq0$,
then the inclusion of $y_{i,t-k}$ eliminates that bias.
I offer no \itx{a priori} view on the strength of $\phi_1$, the autoregressor coefficient.
A significant estimate would imply heavy bias to the static models, validating the use in this dynamic setting.
\par
There are also significant theoretical advantages to introducing the $q$-lags of the spatial spillover.
This specification with both spatial and temporal lags accounts for the possibility
that the spillover effect from firm $j$'s production factor
has a lagged impact on $i$, conceptualisable as the formation of clusters
from firms seeing the benefits of colocation over some time after a productivity shock.
Without including the $q$ lags within $z_{i,t}$,
if $\mathbf z_{t-k}^\text{spill}$ is correlated with $\mathbf z_{t}$,
then estimating $\delta_0$ in a static model when the true effect is dynamic
exposes estimation to omitted variable bias (OVB) by attribution of the lagged to the current period's regressor.
\par
In practice for this dataset however, it is unlikely that more-than-one period spillover lag is relevant.
This is due to this analysis using Fame's annual frequency observations.
As a result, I show the structural equations in an AR(1,0) structure
so that the primary advancement on the static models
are purely the use of the weak exogeneity assumption, and first-order autoregressor.
\begin{align}
    \label{eq:dyn-main-ar10}
    y_{i,t} &= \alpha_i + \gamma_t + \beta_0 + \beta_1 k_{i,t} + \beta_2 l_{i,t} 
                + \phi_1 y_{i,t-k} + \varepsilon_{i,t}
                \\ \notag &\qquad 
                + \delta_k \sum_{j\neq i} w_{ij}y_{j,t}
                + (\xi_k-\delta_k\beta_1) \cdot \sum_{j\neq i} w_{ij}k_{j,t}
                + (\zeta_k-\delta_k\beta_2) \cdot \sum_{j\neq i} w_{ij}l_{j,t}
\end{align}

\subsection{Dynamic panel identification}

The dynamic panel specification uses a autocorrelation error orthogonality assumption
to find valid identification \eqref{eq:dyn-no-ac}.
This replaces \eqref{eq:static-identification-factor-exogeneity} and instead allows
for some moderate level of cross-sectional error correlation \eqref{eq:dyn-xs-corr-errors}.
\begin{align}
    \label{eq:dyn-no-ac}
    \mathbb E[\varepsilon_{i,t}\varepsilon_{i,t-1}] &= 0       ,\qquad \forall t,t-1
    \\
    \label{eq:dyn-xs-corr-errors}
    \mathbb E[\varepsilon_{i,t}\varepsilon_{j,t}] &\neq 0   ,\qquad \exists i,j
    \\
    \label{eq:dyn-error-wn}
    \mathbb E[\varepsilon_t]    =0,
    \mathbb E[\varepsilon_t^2]  =\sigma_t^2,
    \mathbb E[\varepsilon_t|\Omega_{t-1}] &=0                  ,\qquad \forall i
    \qquad\text{(w.n.)}
\end{align}

No autocorrelation of errors may generally be seen as a more credible assumption,
as in a dynamic, correctly-specified model with the correct autoregressive properties,
time autocorrelation as component of the error term are modelled out.
Conceptually, we could also argue that firms make efficient decisions,
so if there is a shock observed in this period, firms ensure that in the next,
any expected factor inputs that are realised by that period fully take advantage of the shock
and its effect is then captured by the input regressors.
Any remaining residual must have been expected and so mean zero (formally a Martingale sequence).
Further econometric analyses could guarantee this assumption by plotting patterns of errors
and check for autocorrelation, which could come from moving average
or omitted latent variable structures to the specification missing here.
This specification can be identified in reduced panel form,
which requires first differences to eliminate fixed effects,
allowing me to retain the use of firm fixed effects.
\par
Estimation under these dynamic panel conditions are developed by \cite{Anderson1981},
with further techniques expanded in \cite{Arellano1991} and \cite{Blundell1998} which improve efficiency.
The estimator given in \cite{Arellano1991} is designed to allow for
endogeneity, limiting only to weak exogeneity as I do here.
Internal instruments of the endogenous regressors are found using their lagged first differences of the log variables.
Identification is improved in the System GMM approach given in \cite{Blundell1998}
by also including the first logged, not differenced, lag of the endogenous regressors
which meets the usual instrument conditions, as an instrument.
I estimate this section's results with
System GMM\footnote{Implemented with \cite{Wu2023}, with heavy modifications in \cite{Lip2026}}.
Valid instruments are therefore theoretically derived as $\Delta x_{i,\tau-j}$ for any $j$
if $x_{i,\tau}$ is a valid instrument.

\subsection{Testing AR structure}
\label{ss:dyn-ar}

I firstly relax the strict exogeneity assumption of the standard factors to weak exogeneity.
There are two approaches to this:
capital and labour can be assumed to be predetermined:
capital faces some time-to-build constraint
where each firm allocates investment in the previous period
based on the limited information available then, $\Omega_{t-1}$.
The predetermined nature of capital is a common modelling choice of
firm production in a dynamic system, though usually with quarterly observations.
A similar sticky decision pattern is assumed of labour,
perhaps through job searching and hiring frictions.
For handling panel data in an autoregressive setting,
first-differencing will lead to the reduced-form error term $\Delta\varepsilon_t$
spanning 2 periods.
As a result, weak exogeneity of the transformed factors
requires lags of $t-2$.
These deep lags are valid instruments which defensible exclusion restriction.
Since the zero autocorrelation property applies to any firm $i$,
the full vector of deep lag factor regressors and their transformations
are plausibly exogenous.
\begin{align}
    \notag
    \text{Pre-determined factors:}
    \\ \notag
    K_{i,t} &= I_{i,t} + dK_{i,t-1}
    \\ \notag
    I_{i,t} &= \mathbb E[I_{i,t}^*|\Omega_{t-1}]
    \\ \notag
    L_{i,t} &= \mathbb E[L_{i,t}^*|\Omega_{t-1}]
    \\ \notag
    K_{i,t}, L_{i,t} &\in \Omega_{t-1}
    \\ \notag
    \mathbb E[\varepsilon_{i,t}|\Omega_{t-1}] &= 0 
    \\
    \label{eq:dyn-weak-exog-pred}
    \mathbb E[
        \varepsilon_{i,t}
        \vert
        \begin{pmatrix}
            k_{i,\tau} & l_{\tau}
        \end{pmatrix}
    ] &= 0_{(1\times t)}
    ,\ \forall i\ \forall \tau\leq t\leq 1
    \\
    \label{eq:dyn-weak-exog-pred-diff}
    \mathbb E[
        \Delta\varepsilon_{i,t}
        \vert
        \begin{pmatrix}
            \Delta \mathbf W^k\mathbf k_{\tau-1} & \Delta \mathbf W^k \mathbf l_{\tau-1}
        \end{pmatrix}
    ] &= 0_{(1\times t)}
    ,\ \forall k\ \forall \tau\leq t\leq 2
\end{align}

Alternatively, capital and labour may be flexible inputs:
at least partially able to react to shocks
as part of the decision set at time $t$.
This may be a more realistic assumption when handling annual data, such as Fame:
shocks occurring at the beginning of a calendar year would have no arguable
reason why the firm's input choices later in that year hadn't reacted to knowledge of the shock,
any more than it had occurred at the end of the previous year.
It is therefore much harder to argue the plausibly exogeneity of contemporaneous factors to the error term.
In practice, as with any dynamic panel model, the procedure laid out in this paper
is not bound by the features of this particular dataset, which happens to be annual.
Future analysis repeating this question with quarterly observations may well
exploit the plausible predetermined nature of some production factors.
\begin{align}
    \notag
    \text{Flexible factors:}
    \\ \notag
    K_{i,t}, L_{i,t} &\in \Omega_{t}
    \\ \notag
    \mathbb E[\varepsilon_{i,t}|\Omega_{t}] &= 0 
    \\
    \label{eq:dyn-weak-exog-flex}
    \mathbb E[
        \varepsilon_{i,t}
        \vert
        \begin{pmatrix}
            k_{i,\tau-1} & l_{\tau-1}
        \end{pmatrix}
    ] &= 0_{(1\times t)}
    ,\ \forall i\ \forall \tau\leq t\leq 1
    \\
    \label{eq:dyn-weak-exog-flex-diff}
    \mathbb E[
        \Delta\varepsilon_{i,t}
        \vert
        \begin{pmatrix}
            \Delta \mathbf W^k\mathbf k_{\tau-2} & \Delta \mathbf W^k \mathbf l_{\tau-2}
        \end{pmatrix}
    ] &= 0_{(1\times t)}
    ,\ \forall k\ \forall \tau\leq t\leq 2
\end{align}

In this dataset, the empirical experiments run suggest that the decision
over whether the classical factors should be predetermined or not
makes no difference for this dataset.
I present estimations of the 
\wraptable
    {Dynamic panel: baseline instrumental specification}
    {../../model/output/results_3a_basic_str}
    {
        \begin{itemize}
        \item[] Hansen J-test for over-identification: `\textsf{X}' indicates insufficient evidence
                at the 95\% level to reject the null-hypothesis that the model is identified.
        \item[] Anderson-Rubin AR(1) test:  `\checkmark' indicates the use of single-order lags is a weak instrument,
                as expected from a first-differencing System GMM approach.
        \item[] Anderson-Rubin AR(2) test: `\textsf{X}' indicates insufficient evidence to
                reject the use of second-order lags due to their being a weak instrument,
                as expected from the first-differencing with lags greater than 2nd-order to overcome contemporaneous endogeneity.
                `checkmark' indicates some instruments used were weak.
        \end{itemize}
    }

Devising an internal instrument requires a similar evaluation of the components of $y_{j,t}$ as in the static problem.
There is a dynamic corollary to the static derivation of the endogenous term $y_{i,t}$.
\begin{align}
    \mathbf y_t
    &= \boldsymbol{\alpha} 
        + \gamma_t \boldsymbol{\iota} 
        + \beta_0 \boldsymbol{\iota} 
        + \beta_1 \mathbf{k}_t 
        + \beta_2 \mathbf{l}_t 
        + \delta_0 \mathbf{W} \mathbf{y}_t
        + (\xi_0 - \delta_0\beta_1) \mathbf{W} \mathbf{k}_t 
        + (\zeta_0 - \delta_0\beta_2) \mathbf{W} \mathbf{l}_t
        + \varepsilon_t \notag \\
    &\quad 
        + \sum_{m=1}^p (\phi_m \mathbb{I} + \delta_m \mathbf{W}) \mathbf{y}_{t-m} 
    % \notag \\ &\quad 
        + \sum_{m=1}^p (\xi_m - \delta_m\beta_1) \mathbf{W} \mathbf{k}_{t-m} 
        + \sum_{m=1}^p (\zeta_m - \delta_m\beta_2) \mathbf{W} \mathbf{l}_{t-m} 
        \\[1.5ex]
    (\mathbb{I} - \delta_0 \mathbf{W})\mathbf{y}_t 
    &= \boldsymbol{\alpha} 
        + \gamma_t \boldsymbol{\iota} 
        + \beta_0 \boldsymbol{\iota} 
        + \beta_1 \mathbf{k}_t 
        + \beta_2 \mathbf{l}_t 
        + (\xi_0 - \delta_0\beta_1) \mathbf{W} \mathbf{k}_t 
        + (\zeta_0 - \delta_0\beta_2) \mathbf{W} \mathbf{l}_t \notag \\
    &\quad 
        + \sum_{m=1}^p (\phi_m \mathbb{I} + \delta_m \mathbf{W}) \mathbf{y}_{t-m} \notag \\
    &\quad 
        + \sum_{m=1}^p \mathbf{W} \Big[ 
            (\xi_m - \delta_m\beta_1)\mathbf{k}_{t-m} 
            + (\zeta_m - \delta_m\beta_2)\mathbf{l}_{t-m} 
        \Big] 
        + \boldsymbol{\varepsilon}_t \\[2ex]
    %
    \mathbf{y}_t 
    &= \sum_{k=0}^\infty \delta_0^k \mathbf{W}^k \Bigg( 
        \boldsymbol{\alpha} 
        + \gamma_t \boldsymbol{\iota} 
        + \beta_0 \boldsymbol{\iota} 
        + \beta_1 \mathbf{k}_t 
        + \beta_2 \mathbf{l}_t 
        + (\xi_0 - \delta_0\beta_1) \mathbf{W} \mathbf{k}_t 
        + (\zeta_0 - \delta_0\beta_2) \mathbf{W} \mathbf{l}_t \notag \\
    &\quad 
        + \sum_{m=1}^p (\phi_m \mathbb{I} + \delta_m \mathbf{W}) \mathbf{y}_{t-m} \notag \\
    &\quad 
        + \sum_{m=1}^p \mathbf{W} \Big[ 
            (\xi_m - \delta_m\beta_1)\mathbf{k}_{t-m} 
            + (\zeta_m - \delta_m\beta_2)\mathbf{l}_{t-m} 
        \Big] 
        + \boldsymbol{\varepsilon}_t 
    \Bigg) \\[2ex]
    %
    y_{j,t} 
    &= \dots 
        + \underbrace{ 
            \sum_{k=0}^\infty \left[ \delta_0^k \mathbf{W}^{k+1} \right]_{(j,)} 
            \begin{pmatrix}\mathbf{k}_t & \mathbf{l}_t\end{pmatrix} 
            \begin{pmatrix} \xi_0 \\ \zeta_0 \end{pmatrix} 
        }_{\text{contemporaneous exogenous}_j} \notag \\
    &\quad 
        + \underbrace{ 
            \sum_{m=1}^p \sum_{k=0}^\infty \left[ \phi_m \delta_0^k \mathbf{W}^k + \delta_m \delta_0^k \mathbf{W}^{k+1} \right]_{(j,)} \mathbf{y}_{t-m} 
        }_{\text{entangled historical inertia \& peer effects}_j} \notag \\
    &\quad 
        + \underbrace{ 
            \sum_{m=1}^p \sum_{k=0}^\infty \left[ \delta_0^k \mathbf{W}^{k+1} \right]_{(j,)} 
            \begin{pmatrix}\mathbf{k}_{t-m} & \mathbf{l}_{t-m}\end{pmatrix} 
            \begin{pmatrix} \xi_m - \delta_m\beta_1 \\ \zeta_m - \delta_m\beta_2 \end{pmatrix} 
        }_{\text{historical exogenous network effects}_j} \notag \\
    &\quad 
        + \underbrace{ 
            \sum_{k=0}^\infty [\delta_0^k \mathbf{W}^k]_{(j,)} \cdot \boldsymbol{\varepsilon}_t 
        }_{\text{error transformation}_j}
\end{align}

For system GMM, I need to show that the reduced-form first difference regressors are exogenous.
The first-order autoregressor, $\Delta ln(y_{it})$ is endogenous.
\begin{align}
    \Delta y_{i,t} &= \Delta \gamma_t
                + \beta_1 \Delta k_{i,t}
                + \beta_2 \Delta l_{i,t}
                + \sum_{k=1}^p \phi_k \Delta y_{i,t-k}
                + \delta_0 \sum_{j\neq i} w_{ij,t} \Delta z_{j,t}
                + \Delta \varepsilon_{i,t}
    \\
    \Delta \varepsilon_{i,t} &= \varepsilon_{i,t} - \varepsilon_{i,t-1}
    \\ \notag
    \Delta y_{i,t-1} &= \qquad\dots\qquad + \Delta\varepsilon_{i,t-1}
    \\ \notag
    &= \qquad\dots\qquad + \varepsilon_{i,t-1} - \varepsilon_{i,t-2}
    \\ \notag
    &\implies \mathbb E[ \Delta y_{i,t-1}\Delta\varepsilon_{i,t}] \neq 0
    \\ \notag
    \Delta y_{i,t-2} &= \qquad\dots\qquad + \varepsilon_{i,t-2} - \varepsilon_{i,t-3}
    \\ \notag
    &\implies \mathbb E[ \Delta y_{i,t-1}\Delta\varepsilon_{i,t}] = 0
\end{align}

This is confirmed by the Anderson-Rubin (AR(1)) test in my results,
which rejects the null hypothesis at the 1\% level for all models,
suggesting that first-order lags of log GVA are a weak instrument
\todo{Why? Shouldn't they fail exclusion restriction, not instrument correlation}.
In addition, I assume that capital, $k_{i,t}$ is pre-determined, as investment is chosen
one period ahead, in line with standard Perpetual Inventory Method approaches\footnote{
    Noting that this is ordinarily done on a quarterly basis,
    however, due to data limitations we adopt the same thing here even though Fame
    provides observations annually.
}\todo{Review predetermined material. Additionally, explain the logic of instrumenting for capital}
\par
To overcome this endogeneity, I introduce deeper autoregressive lags as instrumental variables,
as proposed by \cite{Anderson1981}.
I use 3 autoregressive instrument lags starting from $t-2$:
$\ln(y_{i,t-2}), \ln(y_{i,t-2}), \ln(y_{i,t-p-3})$ for log GVA,
and 2 lags for capital $\ln(k_{i,t-p-1}), ln(k_{t-p-2})$.
An Anderson-Rubin (AR(2)) test shows that deep lags are a valid instrument for log GVA.
The test applied to the second lag of $ln(gva_1)$
fails to reject the null hypothesis of for any specification,
so I assume it is a valid instrument.
\par
Additionally, using multiple deep lags allows for overidentification:
in this basic AR(1) specification, I have $k=4$ regressors and $l=7$ instruments.
I could increase this. Fame is rich on the time dimension, with up to 19 lags at most for a firm,
and a median of 6 lags, so I could safely provide more lags of regressors as instruments.
Further lags is likely to have a decreasing marginal effect on explanatory power
and additionally \todo(...)
\par
I estimate the equation fulfilling the vector of moment conditions,
(\ref{eq:dyn-moment-conditions}) via GMM.
\begin{align}
    \label{eq:dyn-moment-conditions}
    \mathbb E \left[
        \Delta\varepsilon_{it}
        \cdot
        \begin{pmatrix}
            \Delta \ln(y_{i,t-2})     \\
            \Delta \ln(y_{i,t-3})     \\
            \Delta \ln(y_{i,t-4})     \\
            \Delta \ln(k_{i,t-2})     \\
            \Delta \ln(k_{i,t-3})     \\
            \Delta \Sigma_jw_{ijt}z_{jt}
        \end{pmatrix}
    \right]
    &= 0
\end{align}

\wraptable
    {Dynamic panel: deeper instrumental lags}
    {../../model/output/results_3b_dyn_instruments_str}
    {
        \begin{itemize}
        \item[] Hansen J-test for over-identification: `\textsf{X}' indicates insufficient evidence
                at the 95\% level to reject the null-hypothesis that the model is identified.
        \item[] Anderson-Rubin AR(1) test:  `\checkmark' indicates the use of single-order lags is a weak instrument,
                as expected from a first-differencing System GMM approach.
        \item[] Anderson-Rubin AR(2) test: `\textsf{X}' indicates insufficient evidence to
                reject the use of second-order lags due to their being a weak instrument,
                as expected from the first-differencing with lags greater than 2nd-order to overcome contemporaneous endogeneity.
                `checkmark' indicates some instruments used were weak.
        \end{itemize}
    }


% \wraptable
%     {Dynamic panel: deeper vector regressor lags}
%     {../../model/output/instr_3b_Jover}
%     {Hansen J-test for over-identification: `\textsf{X}' indicates insufficient evidence
%     at the 95\% level to reject the null-hypothesis that the model is identified.
%     Anderson-Rubin AR(1) test: 
%         Model includes a constant as the baseline for \itx{t} fixed effects.
%     }

% \wraptable
%     {Dynamic panel: deeper autoregressive lags}
%     {../../model/output/instr_3c_ar}
%     {
%         Model includes a constant as the baseline for \itx{t} fixed effects.
%     }

\par 
Finally, a Sargan-Hansen J-test is implemented to check if these instruments are
over-identified, meaning the moment conditions cannot be fulfilled with 95\% confidence.
In all models, the J-test fails to reject the null hypothesis,
meaning the model's moment conditions are jointly correct.

\todo{!!!!}
Following above reasoning, lags of $\Delta\sum_{j\neq i} w_{i,j,t}z_{j,t}$
of a higher order than 2 are exogenous to my reduced-form specification.
However, 